\section{Design of Apollo-RTOS}

\subsection{Inspiration}

Apollo-RTOS draws inspiration from both the Apollo Guidance Computer (AGC) operating
system and the modern Linux OS. The AGC OS was designed to meet specific mission
requirements, allowing for a simplified and optimized design. In contrast, Linux is a
general-purpose OS built to support diverse applications, resulting in a much more
complex architecture.

The AGC operated in a tightly controlled hardware and software environment where all
components could be trusted. Thus, isolating the hardware, OS, and application code
was unnecessary. Additionally, the AGC lacked memory protection features, making
kernel/user space isolation---a fundamental aspect of general-purpose OSs like
Linux---impossible.

Apollo-RTOS adopts the same high-level design philosophy as the AGC: implement core
concepts effectively in a resource-constrained environment, while applying
contemporary engineering where beneficial. Key design decisions include:

\begin{description}[leftmargin=15pt, labelindent=0pt]
    \item [Hardware Selection] The \texttt{micro:bit} v1 was chosen as the primary
        target platform. Its Cortex-M0 CPU is among the smallest in the ARM family and
        closely resembles the AGC hardware. 
        \ifextraC
            It also includes integrated peripherals, providing practical examples to
            demonstrate OS functionality.
        \fi

    \item [Monolithic Kernel Design] Following Linux's model, Apollo-RTOS adopts a
        monolithic kernel for simplicity and efficient execution in a constrained
        environment, at the cost of more challenging verification and testing.

    \item [No Kernel/User Space Isolation] In microcontroller programming, the entire
        program---kernel and applications---is compiled into a single binary. Strict
        isolation would add unnecessary complexity.

    \item [Library Calls Instead of System Calls] Since Apollo-RTOS does not enforce
        user/kernel space separation, direct library calls are used rather than
        system calls, simplifying the interface and reducing overhead.

    \item [Shared Memory IPC] Although shared memory risks race conditions, it offers
        the fastest and simplest form of IPC and is thus adopted as the primary
        mechanism.

    \item [Memory Usage Philosophy] While the \texttt{micro:bit} has limited memory,
        it offers far more than the AGC. Prioritizing maintainability over aggressive
        optimization, we follow Knuth’s maxim that ``premature optimization is the root
        of all evil.''
\end{description}

We also drew inspiration from the \texttt{microbian} OS for the \texttt{micro:bit},
developed by \citet{microbian}, particularly for hardware headers, startup code, and
interrupt handler structures.

\ifextraA
For the language and development ecosystem, we chose modern C++ over C, due to:
\begin{itemize}
    \item Concise syntax (classes, templates, RAII) that reduces boilerplate.
    \item Greater optimization opportunities through templates, \texttt{constexpr},
        and references.
    \item Clean, expressive abstractions via type deduction, improving maintainability.
\end{itemize}
\fi

\subsection{Implementation Architecture}

In this section, we briefly overview Apollo-RTOS's core modules and their
interconnections:

\begin{description}[leftmargin=2em, labelindent=0pt]
    \item [Scheduler] The centerpiece of the OS, managing:
        \begin{itemize}[left=0pt]
            \item Process execution and context switching.
            \item Library functions for process creation, priority management, and
                synchronization.
        \end{itemize}

    \item[Waitlist] Implements a deadline-driven preemptive scheduler (similar to cron
        jobs), executing time-sensitive tasks at precise intervals.

    \item[Signals, Mutexes, and Barriers] Enable low-overhead inter-process
        communication (IPC) and synchronization among processes.

    \item[Memory Allocator] Handles all dynamic memory management. Allocators maintain
        freelists of memory chunks and serve both runtime heap and stack allocations.

    \item[File System] Provides persistent storage through:
        \begin{itemize}[left=0pt]
            \item A FRAM-based file system over I2C.
            \item An optional flash memory file system.
            \item Support for both memory-mapped and character-stream file access.
        \end{itemize}

    \item[Error Handling] Provides diagnostic logging and error detection, including
        runtime assertions and fault recovery.

    \item[Recovery] Supports wrapping critical sections and processes with automatic
        recovery against different types of resets.

    \item[Drivers] Interface with hardware components:
        \begin{itemize}[left=0pt]
            \item \textbf{Timer}: Two 16-bit timers are configured for system clocking
                and profiling (\lstinline[style=myagc]|TIMER1| and
                \lstinline[style=myagc]|TIMER2|).
            \item \textbf{Serial}: Manages UART-based serial communication.
            \item \textbf{I2C}: Manages I2C master communication.
            \item \textbf{Display}: Controls the \(5 \times 5\) LED matrix.
        \end{itemize}

    \item[Shell] Provides a Unix-like command-line interface for user interaction and
        process management.

    \item[Testing Framework] Allows registering and running unit and subsystem tests in
        isolation during system startup.

    \item[``Standard" Library] Offers utility data structures and functions, with
        interfaces inspired by the C++ Standard Library.
\end{description}
